% Yêu cầu 1 — Ngân sách: ~4 trang
\chapter{Kiến trúc và nguyên lý hoạt động}\label{chap:kientruc}

\section{Tổng quan về tường lửa ứng dụng web}
Tường lửa ứng dụng web (WAF) là thành phần kiểm tra nội dung giao dịch HTTP/HTTPS giữa người dùng và ứng dụng web, nhằm phát hiện và ngăn chặn các tấn công ở tầng ứng dụng như SQL Injection, Cross-Site Scripting hay Local File Inclusion~\cite{owasptop10}. Khác với tường lửa mạng chỉ xét địa chỉ IP, cổng và trạng thái kết nối, WAF phân tích cú pháp yêu cầu HTTP: đường dẫn, tham số, header, cookie và thân yêu cầu.

WAF có thể được triển khai theo ba mô hình: \textbf{reverse proxy} (WAF đứng trước ứng dụng, kết thúc TLS, kiểm tra rồi mới chuyển tiếp yêu cầu; đây là mô hình nhóm dùng trong lab), \textbf{nhúng} (chạy như một module trong máy chủ web Apache, Nginx) và \textbf{ngoài luồng} (nhận bản sao lưu lượng để phát hiện, không chặn trực tiếp).

Về chính sách, WAF có thể theo \textbf{mô hình âm} (negative security model: chặn những gì khớp mẫu tấn công đã biết) hoặc \textbf{mô hình dương} (positive security model: chỉ cho phép những gì đã được khai báo hợp lệ). Bộ OWASP CRS chủ yếu theo mô hình âm, với các luật tổng quát không gắn với một ứng dụng cụ thể.

\section{ModSecurity}
\subsection{Lịch sử và kiến trúc}
ModSecurity do Ivan Ristić phát triển từ năm 2002, ban đầu là một module của Apache HTTP Server~\cite{modsecbook}. Dự án sau đó được Breach Security rồi Trustwave SpiderLabs bảo trợ; năm 2024, khi Trustwave kết thúc hỗ trợ, dự án được chuyển giao cho OWASP~\cite{modsecurity}.

Có hai dòng phiên bản chính:
\begin{itemize}
  \item \textbf{ModSecurity v2:} gắn chặt với Apache, phụ thuộc vào cấu trúc dữ liệu nội bộ của Apache.
  \item \textbf{ModSecurity v3 (libmodsecurity):} viết lại bằng C++ thành một thư viện độc lập với máy chủ web. Máy chủ web tích hợp thông qua \emph{connector}, phổ biến nhất là \texttt{ModSecurity-nginx}. Mỗi yêu cầu HTTP được thư viện xử lý như một \emph{transaction}: connector chuyển từng phần dữ liệu (header, body) vào thư viện và nhận lại quyết định cho qua hay chặn.
\end{itemize}

ModSecurity bản thân chỉ là một \emph{bộ máy thực thi luật}; khả năng phát hiện tấn công đến từ bộ luật nạp vào, mà phổ biến nhất là OWASP CRS (mục~\ref{sec:crs}).

\subsection{Ngôn ngữ luật SecLang}
Luật và cấu hình ModSecurity được viết bằng ngôn ngữ SecLang. Một luật có cấu trúc:
\begin{center}
\texttt{SecRule BIẾN TOÁN\_TỬ "HÀNH\_ĐỘNG"}
\end{center}
\begin{itemize}
  \item \textbf{Biến} xác định phần dữ liệu cần kiểm tra, ví dụ \texttt{ARGS} (tham số), \texttt{REQUEST\_HEADERS}, \texttt{REQUEST\_URI}, \texttt{REQUEST\_BODY}, hay biến giao dịch \texttt{TX}.
  \item \textbf{Toán tử} xác định cách so khớp: \texttt{@rx} (biểu thức chính quy), \texttt{@pm} (khớp nhiều chuỗi), \texttt{@detectSQLi} và \texttt{@detectXSS} (dùng thư viện libinjection~\cite{libinjection} để phân tích cú pháp thay vì so khớp mẫu).
  \item \textbf{Hành động} gồm: định danh (\texttt{id}, \texttt{phase}, \texttt{msg}, \texttt{severity}, \texttt{tag}); biến đổi dữ liệu trước khi so khớp (\texttt{t:urlDecodeUni}, \texttt{t:lowercase}, \texttt{t:htmlEntityDecode}) để chống né tránh bằng mã hóa; hành động ngắt (\texttt{deny}, \texttt{block}, \texttt{pass}, \texttt{drop}); và gán biến (\texttt{setvar}).
\end{itemize}

\begin{lstlisting}[caption={Một luật SecLang phát hiện SQLi bằng libinjection},label={lst:seclang}]
SecRule ARGS "@detectSQLi" \
    "id:100001,phase:2,block,t:none,t:urlDecodeUni,\
    log,msg:'SQL Injection (libinjection)',severity:'CRITICAL',\
    setvar:'tx.inbound_anomaly_score_pl1=+%{tx.critical_anomaly_score}'"
\end{lstlisting}

Luật trong Đoạn mã~\ref{lst:seclang} kiểm tra mọi tham số ở pha 2, giải mã URL trước khi phân tích, và nếu phát hiện SQLi thì ghi log và cộng điểm bất thường thay vì chặn ngay. Cách cộng điểm này là nền tảng của cơ chế anomaly scoring trong CRS.

Bên cạnh luật, một số chỉ thị cấu hình quan trọng:
\begin{itemize}
  \item \texttt{SecRuleEngine On | Off | DetectionOnly}: bật chặn, tắt, hoặc chỉ phát hiện.
  \item \texttt{SecRequestBodyAccess} / \texttt{SecResponseBodyAccess}: cho phép kiểm tra thân yêu cầu / phản hồi.
  \item \texttt{SecAuditEngine}, \texttt{SecAuditLogParts}, \texttt{SecAuditLogFormat}: điều khiển nhật ký kiểm toán (audit log), có thể xuất dạng JSON để đưa vào hệ thống SIEM.
\end{itemize}

\subsection{Các pha xử lý giao dịch HTTP}
ModSecurity chia mỗi giao dịch thành năm pha (Hình~\ref{fig:phases}); mỗi luật được gán vào một pha qua hành động \texttt{phase}:
\begin{enumerate}
  \item \textbf{Request headers:} ngay sau khi nhận xong header yêu cầu. Phù hợp với luật về phương thức, URI, header.
  \item \textbf{Request body:} sau khi đọc thân yêu cầu (form, JSON, multipart). Phần lớn luật phát hiện tấn công chạy ở pha này vì đã có đủ tham số.
  \item \textbf{Response headers:} trước khi gửi header phản hồi về người dùng.
  \item \textbf{Response body:} kiểm tra nội dung phản hồi, ví dụ phát hiện rò rỉ thông báo lỗi SQL hay mã nguồn.
  \item \textbf{Logging:} sau khi giao dịch kết thúc; chỉ ghi log, không thể chặn.
\end{enumerate}

\begin{figure}[!htb]
\centering
\begin{tikzpicture}[
  box/.style={draw,rounded corners,minimum width=2.3cm,minimum height=9mm,font=\footnotesize,align=center},
  end/.style={draw,minimum width=1.7cm,minimum height=9mm,font=\footnotesize,align=center,fill=gray!15},
  lbl/.style={font=\scriptsize}]
  \node[end] (cli) at (0,0) {Người\\dùng};
  \node[box] (p1) at (3,0) {Pha 1\\Request headers};
  \node[box] (p2) at (6.2,0) {Pha 2\\Request body};
  \node[end] (app) at (9.4,0) {Ứng dụng\\web};
  \node[box] (p3) at (6.2,-2) {Pha 3\\Response headers};
  \node[box] (p4) at (3,-2) {Pha 4\\Response body};
  \node[draw,rounded corners,fill=gray!8,font=\footnotesize] (p5) at (4.6,-3.4) {Pha 5 -- Logging: ghi audit log sau khi giao dịch kết thúc};
  \draw[-Stealth] (cli) -- (p1);
  \draw[-Stealth] (p1) -- (p2);
  \draw[-Stealth] (p2) -- node[lbl,above]{cho qua} (app);
  \draw[-Stealth] (app) |- (p3);
  \draw[-Stealth] (p3) -- (p4);
  \draw[-Stealth] (p4) -| (cli);
  \draw[-Stealth,dashed,red!70!black] (p2.north) -- ++(0,0.6) -| (cli.north);
  \node[lbl,text=red!70!black] at (3,0.85) {điểm inbound $\geq$ ngưỡng: chặn, trả 403};
\end{tikzpicture}
\caption{Các pha xử lý của ModSecurity/Coraza khi đặt làm reverse proxy}
\label{fig:phases}
\end{figure}

\section{Coraza}
OWASP Coraza là một WAF mã nguồn mở viết bằng Go, cài đặt lại bộ máy SecLang của ModSecurity~\cite{coraza}. Coraza dùng cùng cú pháp luật và chạy được nguyên bộ OWASP CRS, nên chính sách viết cho ModSecurity có thể chuyển sang Coraza mà gần như không phải sửa. Khác biệt chính:
\begin{itemize}
  \item \textbf{Ngôn ngữ và an toàn bộ nhớ:} Go loại bỏ nhiều lớp lỗi bộ nhớ thường gặp ở mã C/C++.
  \item \textbf{Tích hợp:} Coraza là thư viện Go, nhúng trực tiếp vào ứng dụng Go hoặc dùng qua các tích hợp với Caddy, Envoy/Istio (proxy-wasm), HAProxy (SPOA), thay vì connector cho Nginx/Apache.
  \item \textbf{Tình trạng dự án:} Coraza được phát triển tích cực trong hệ sinh thái cloud-native, trong khi ModSecurity v3 hiện chủ yếu được sửa lỗi và vá bảo mật.
\end{itemize}
Vì hai phần mềm dùng chung SecLang và CRS, các nguyên lý trình bày ở mục này và mục~\ref{sec:crs} áp dụng cho cả hai. Lab của nhóm dùng ModSecurity v3 trên Nginx.

\section{OWASP Core Rule Set}\label{sec:crs}
\subsection{Cấu trúc bộ luật}
OWASP Core Rule Set (CRS) là bộ luật phát hiện tấn công tổng quát dùng cho ModSecurity, Coraza và nhiều WAF tương thích~\cite{crs}. Các luật được chia thành tệp theo nhóm tấn công, đánh số theo dải ID; Bảng~\ref{tab:crsfiles} liệt kê các nhóm liên quan đến kịch bản demo. Cấu hình chung (ngưỡng điểm, paranoia level) nằm trong \texttt{crs-setup.conf}.

\begin{table}[!htb]
\centering
\caption{Một số nhóm luật CRS liên quan đến đề tài}
\label{tab:crsfiles}
\small
\begin{tabularx}{\linewidth}{@{}llX@{}}
\toprule
\textbf{Dải ID} & \textbf{Tệp luật} & \textbf{Nội dung} \\ \midrule
920xxx & REQUEST-920-PROTOCOL-ENFORCEMENT & Kiểm tra tuân thủ giao thức HTTP \\
930xxx & REQUEST-930-APPLICATION-ATTACK-LFI & Path traversal, truy cập tệp hệ thống \\
932xxx & REQUEST-932-APPLICATION-ATTACK-RCE & Chèn lệnh hệ điều hành \\
941xxx & REQUEST-941-APPLICATION-ATTACK-XSS & Cross-Site Scripting \\
942xxx & REQUEST-942-APPLICATION-ATTACK-SQLI & SQL Injection \\
949xxx & REQUEST-949-BLOCKING-EVALUATION & So sánh điểm inbound với ngưỡng \\
959xxx & RESPONSE-959-BLOCKING-EVALUATION & So sánh điểm outbound với ngưỡng \\
\bottomrule
\end{tabularx}
\end{table}

\subsection{Cơ chế chấm điểm bất thường (anomaly scoring)}
CRS mặc định hoạt động theo cơ chế chấm điểm. Mỗi luật khi khớp không chặn ngay mà cộng vào điểm bất thường của giao dịch một giá trị theo mức nghiêm trọng: CRITICAL cộng 5, ERROR cộng 4, WARNING cộng 3, NOTICE cộng 2. Cuối pha 2, luật trong tệp 949 so tổng điểm inbound với ngưỡng (mặc định 5); nếu đạt hoặc vượt ngưỡng, yêu cầu bị chặn và audit log ghi luật 949110 kèm danh sách các luật đã góp điểm. Tương tự, cuối pha 4 tệp 959 đánh giá điểm outbound (ngưỡng mặc định 4).


Cơ chế này có hai ưu điểm. Thứ nhất, một yêu cầu chỉ khớp một luật nhẹ (ví dụ NOTICE) sẽ không bị chặn, giảm chặn nhầm. Thứ hai, quản trị viên có thể điều chỉnh độ nhạy bằng một tham số duy nhất là ngưỡng, và đọc được trong log vì sao một yêu cầu bị chặn.

\subsection{Paranoia level}
Mỗi luật CRS được gắn một mức paranoia (PL) từ 1 đến 4:
\begin{itemize}
  \item \textbf{PL1} (mặc định): luật chắc chắn, rất ít chặn nhầm; phù hợp hầu hết ứng dụng.
  \item \textbf{PL2}: thêm luật bắt các biến thể tấn công khó hơn; bắt đầu cần tinh chỉnh.
  \item \textbf{PL3, PL4}: luật rất chặt (ví dụ hạn chế ký tự đặc biệt), phát hiện tốt hơn nhưng chặn nhầm nhiều, đòi hỏi viết luật loại trừ cho từng ứng dụng.
\end{itemize}
CRS~4 dùng biến \texttt{tx.blocking\_paranoia\_level} để chọn mức luật được phép góp điểm chặn, và \texttt{tx.detection\_paranoia\_level} để chạy thêm các luật ở mức cao hơn nhưng chỉ ghi log. Nhờ đó có thể quan sát tác động của một mức cao hơn trên lưu lượng thật trước khi bật chặn.

\subsection{Chế độ DetectionOnly và Blocking}
Chỉ thị \texttt{SecRuleEngine} quyết định WAF có thực thi hành động chặn hay không:
\begin{itemize}
  \item \textbf{DetectionOnly:} mọi luật vẫn được đánh giá và ghi log, nhưng hành động ngắt (\texttt{deny}, \texttt{block}) không được thực thi. Quy trình khuyến nghị là chạy chế độ này với PL1 khi mới triển khai, thu thập các trường hợp chặn nhầm và viết luật loại trừ, rồi mới chuyển sang Blocking.
  \item \textbf{On (Blocking):} yêu cầu vượt ngưỡng bị chặn, mặc định trả về mã 403.
\end{itemize}
